\documentclass[UTF8, fontset=fandol,12pt, a4paper]{ctexart}
\usepackage{amsmath, amssymb}
\usepackage[margin=2.5cm]{geometry}
\usepackage{pgfplots}
\pgfplotsset{compat=1.18}
\usepackage{titling}
\setlength{\droptitle}{-1.5cm}
\setlength{\parskip}{0.4em}
\ctexset{section={format=\large\bfseries,beforeskip=1.2em,afterskip=0.6em}}

\title{\LARGE 信号与系统分析\\[0.3em]\large 第二次作业}
\author{黄泱博 \quad 2024012867}
\date{\today}

\begin{document}
\pagestyle{plain}
\maketitle

\section* {2.3 求卷积并绘图}

输入与冲激响应分别为
\[
  x[n]=\left(\frac12\right)^{n-2}u[n-2],
  \qquad h[n]=u[n+2].
\]
离散时间卷积为
\begin{align*}
  y[n] &= x[n]*h[n]
       = \sum_{k=-\infty}^{\infty}x[k]h[n-k] \\
       &= \sum_{k=-\infty}^{\infty}
          \left(\frac12\right)^{k-2}u[k-2]u[n-k+2].
\end{align*}
要求 $2\le k\le n+2$，只有当 $n\ge0$ 时求和区间非空。
\[
  y[n]=\begin{cases}
      0, & n<0,\\[0.4em]
      \displaystyle\sum_{k=2}^{n+2}\left(\frac12\right)^{k-2}
      =\frac{1-(1/2)^{n+1}}{1-1/2}, & n\ge0.
    \end{cases}
\]
最终结果为
\[
  \boxed{y[n]=\left[2-\left(\frac12\right)^n\right]u[n].}
\]

\begin{center}
\begin{tikzpicture}
  \begin{axis}[
    width=0.85\textwidth,height=5.4cm,
    axis lines=middle,
    xlabel={$n$},ylabel={$y[n]$},
    xmin=-3.6,xmax=9.8,ymin=-0.2,ymax=2.45,
    xtick={-3,-2,...,9},ytick={1,1.5,2},
    tick label style={font=\small},
  ]
    \addplot[gray,dashed,domain=-3.3:9.3] {2};
    \addplot[ycomb,blue!70!black,thick,mark=*,mark size=2pt]
      coordinates {(-3,0) (-2,0) (-1,0)};
    \addplot[ycomb,blue!70!black,thick,mark=*,mark size=2pt,
      domain=0:9,samples=10] {2-(1/2)^x};
  \end{axis}
\end{tikzpicture}
\end{center}

\section* {2.8 求卷积并绘图}

输入与冲激响应分别为
\[
  x(t)=\begin{cases}
    t+1, &0 \le t \le 1,\\[0.3em]
    2-t, &1<t \le 2,\\[0.3em]
    0, &\text{其他},
  \end{cases}
  \qquad h(t)=\delta(t+2)+2\delta(t+1).
\]

\begin{align*}
  y(t) &= x(t)*h(t) \\
       &= x(t)*\delta(t+2)+2x(t)*\delta(t+1) \\
       &= x(t+2)+2x(t+1).
\end{align*}

最终结果为
\[
  \boxed{
  y(t)=\begin{cases}
    t+3, &-2\le t<-1,\\[0.3em]
    4, &t=-1,\\[0.3em]
    t+4, &-1<t\le0,\\[0.3em]
    2-2t, &0<t\le1,\\[0.3em]
    0, &\text{其他}.
  \end{cases}}
\]

\begin{center}
\begin{tikzpicture}
  \begin{axis}[
    width=0.85\textwidth,height=6cm,
    axis lines=middle,
    xlabel={$t$},ylabel={$y(t)$},
    xmin=-3,xmax=2,ymin=-0.45,ymax=4.65,
    xtick={-2,-1,0,1},ytick={1,2,3,4},
    tick label style={font=\small},
  ]
    \addplot[blue!70!black,thick] coordinates {(-2.8,0) (-2,0)};
    \addplot[blue!70!black,thick] coordinates {(-2,1) (-1,2)};
    \addplot[blue!70!black,thick] coordinates {(-1,3) (0,4)};
    \addplot[blue!70!black,thick] coordinates {(0,2) (1,0) (1.8,0)};
    \addplot[only marks,mark=*,mark size=2.5pt,
      blue!70!black,mark options={fill=white}]
      coordinates {(-2,0) (-1,2) (-1,3) (0,2)};
    \addplot[only marks,mark=*,mark size=2.5pt,blue!70!black]
      coordinates {(-2,1) (-1,4) (0,4) (1,0)};
  \end{axis}
\end{tikzpicture}
\end{center}

\section* {2.16 判断正误}
(a)正确 
(b)错误
(c)正确
(d)正确

\section* {2.22 求响应并画简图}

\subsection* {(b)}
\begin{align*}
  y(t) &= x(t)*h(t)
        = \int_{-\infty}^{\infty}h(\tau)x(t-\tau)\,d\tau \\
       &= \frac{4}{3}\int_{0}^{1}x(t-\tau)\,d\tau
          -\frac{1}{3}x(t-2) \\
       &= \frac{4}{3}\int_{0}^{1}\bigl[a(t-\tau)+b\bigr]\,d\tau
          -\frac{1}{3}\bigl[a(t-2)+b\bigr] \\
       &= \boxed{at+b}.
\end{align*}

\begin{center}
\begin{tikzpicture}[x=1.6cm,y=1.1cm]
  \draw[->] (-2.3,0) -- (3.4,0) node[right] {$t$};
  \draw[->] (0,-0.9) -- (0,3.8) node[above] {$y(t)$};
  \draw[blue!70!black,thick] (-2,-0.5) -- (3,3.25)
    node[above left,black] {$y(t)=at+b$};
  \fill[blue!70!black] (0,1) circle (1.8pt);
  \node[left] at (0,1) {$b$};
  \node[below left] at (0,0) {$0$};
\end{tikzpicture}

\end{center}

\newpage
\subsection* {(e)}
由于 $x(t)$ 的周期性，只需考虑
\[
  t^*=t-2k\in\left[-\frac12,\frac32\right],\qquad k\in\mathbb{Z}.
\]
\begin{align*}
  y(t) &= x(t)*h(t)
        = \int_{-\infty}^{\infty}x(t-\tau)h(\tau)\,d\tau \\
       &= \int_{0}^{1}x(\tau-t)(1-\tau)\,d\tau \\
       &= \begin{cases}
          \begin{aligned}
            &\displaystyle\int_{0}^{\frac12-t^*}(1-\tau)\,d\tau
             -\int_{\frac12-t^*}^{1}(1-\tau)\,d\tau\\[0.3em]
            &\displaystyle=-(t^*)^2-t^*+\frac14,
          \end{aligned}
          &-\dfrac12<t^*<\dfrac12,\\[1.2em]
          \begin{aligned}
            &\displaystyle-\int_{0}^{\frac32-t^*}(1-\tau)\,d\tau
             +\int_{\frac32-t^*}^{1}(1-\tau)\,d\tau\\[0.3em]
            &\displaystyle=(t^*)^2-t^*-\frac14,
          \end{aligned}
          &\dfrac12<t^*<\dfrac32.
        \end{cases}
\end{align*}

\begin{center}
\begin{tikzpicture}
  \begin{axis}[
    width=0.9\textwidth,height=5.8cm,
    axis lines=middle,
    xlabel={$t$},ylabel={$y(t)$},
    xmin=-2.8,xmax=3.85,ymin=-0.7,ymax=0.75,
    xtick={-2.5,-1.5,-0.5,0.5,1.5,2.5,3.5},
    xticklabels={$-\frac52$,$-\frac32$,$-\frac12$,$\frac12$,
      $\frac32$,$\frac52$,$\frac72$},
    ytick={-0.5,0.5},yticklabels={$-\frac12$,$\frac12$},
    tick label style={font=\small},
  ]
    \foreach \k in {-1,0,1} {
      \addplot[blue!70!black,thick,samples=80,
        domain={2*\k-0.5}:{2*\k+0.5}]
        {-(x-2*\k)^2-(x-2*\k)+0.25};
      \addplot[blue!70!black,thick,samples=80,
        domain={2*\k+0.5}:{2*\k+1.5}]
        {(x-2*\k)^2-(x-2*\k)-0.25};
    }
    \addplot[only marks,mark=*,mark size=2.2pt,
      blue!70!black,mark options={fill=blue!70!black}]
      coordinates {(-2.5,0.5) (-1.5,-0.5) (-0.5,0.5)
        (0.5,-0.5) (1.5,0.5) (2.5,-0.5) (3.5,0.5)};
  \end{axis}
\end{tikzpicture}

{\small 按上述分段式绘制，周期为 $2$。}
\end{center}

\section* {2.24 已知级联及$h_2[n]$}

\subsection* {(a)}
易知：
\[
  u[n-n_1]*u[n-n_2]=(n-n_1-n_2+1)u[n-n_1-n_2],
\]
有：
\begin{align*}
  h_2[n]*h_2[n]
    &=\bigl(u[n]-u[n-2]\bigr)*\bigl(u[n]-u[n-2]\bigr) \\
    &=(n+1)u[n]-2(n-1)u[n-2]+(n-3)u[n-4] \\
    &=\begin{cases}
        1, &n=0,2,\\[0.3em]
        2, &n=1,\\[0.3em]
        0, &\text{其他}.
      \end{cases}
\end{align*}
又由：
\begin{align*}
  h[n] &=h_2[n]*h_2[n]*h_1[n] \\
       &=\bigl(h_2[n]*h_2[n]\bigr)*h_1[n] \\
       &=h_1[n]+2h_1[n-1]+h_1[n-2].
\end{align*}

解方程得：
\[
  \boxed{
  h_1[n]=\begin{cases}
    1, &n=0,4,\\[0.3em]
    3, &n=1,2,\\[0.3em]
    2, &n=3,\\[0.3em]
    0, &\text{其他}.
  \end{cases}}
\]

\newpage
\subsection* {(b)}
由于：
\[
  x[n]=\delta[n]-\delta[n-1],
\]
有：
\begin{align*}
  y[n] &=h[n]*x[n]=h[n]-h[n-1] \\
       &=\boxed{\begin{cases}
          1, &n=0,3,\\[0.3em]
          4, &n=1,\\[0.3em]
          5, &n=2,\\[0.3em]
          -3, &n=4,6,\\[0.3em]
          -4, &n=5,\\[0.3em]
          -1, &n=7,\\[0.3em]
          0, &\text{其他}.
        \end{cases}}
\end{align*}

\section* {2.31 已知差分方程求响应}

输入为
\[
  x[n]=\begin{cases}
    1, &n=-2,3,\\[0.3em]
    2, &n=-1,1,2,\\[0.3em]
    3, &n=0,\\[0.3em]
    0, &\text{其他}.
  \end{cases}
\]
差分方程为
\[
  y[n]+2y[n-1]=x[n]+2x[n-2].
\]

由初始松弛条件有 $y[n]=0,\ \forall n<-2$。逐项代入得
\begin{align*}
  y[-2]+2y[-3] &=x[-2]+2x[-4]=1
    &&\Longrightarrow y[-2]=1,\\
  y[-1]+2y[-2] &=x[-1]+2x[-3]=2
    &&\Longrightarrow y[-1]=0,\\
  y[0]+2y[-1] &=x[0]+2x[-2]=5
    &&\Longrightarrow y[0]=5,\\
  y[1]+2y[0] &=x[1]+2x[-1]=6
    &&\Longrightarrow y[1]=-4,\\
  y[2]+2y[1] &=x[2]+2x[0]=8
    &&\Longrightarrow y[2]=16,\\
  y[3]+2y[2] &=x[3]+2x[1]=5
    &&\Longrightarrow y[3]=-27,\\
  y[4]+2y[3] &=x[4]+2x[2]=4
    &&\Longrightarrow y[4]=58,\\
  y[5]+2y[4] &=x[5]+2x[3]=2
    &&\Longrightarrow y[5]=-114.
\end{align*}

当 $n\ge6$ 时，$x[n]=x[n-2]=0$，故
\[
  y[n]=-2y[n-1],\qquad
  y[n]=y[5](-2)^{n-5}=-114(-2)^{n-5}.
\]
最终结果为
\[
  \boxed{
  y[n]=\begin{cases}
    0, &n<-2\ \text{或}\ n=-1,\\[0.3em]
    1, &n=-2,\\[0.3em]
    5, &n=0,\\[0.3em]
    -4, &n=1,\\[0.3em]
    16, &n=2,\\[0.3em]
    -27, &n=3,\\[0.3em]
    58, &n=4,\\[0.3em]
    -114(-2)^{n-5}, &n\ge5.
  \end{cases}}
\]

\section* {2.38 画出差分因果LTI的框图表示}

\subsection *{(a) $y[n]=\frac{1}{3}y[n-1]+\frac{1}{2}x[n]$}

\begin{center}
\begin{tikzpicture}[
  block/.style={draw,minimum width=1.1cm,minimum height=0.8cm},
  sum/.style={draw,circle,minimum size=0.65cm,inner sep=0pt},
  >=stealth,thick]
  \node[block] (inputgain) at (-2,0) {$\frac12\delta[n]$};
  \node[sum] (sum) at (0,0) {$+$};
  \node[block] (feedbackgain) at (1.25,-2) {$\frac13\delta[n-1]$};
  \draw[->] (-4,0) node[left] {$x[n]$} -- (inputgain.west);
  \draw[->] (inputgain.east) -- (sum.west);
  \draw[->] (sum.east) -- (4,0) node[right] {$y[n]$};
  \fill (2.5,0) circle (1.8pt);
  \draw[->] (2.5,0) |- (feedbackgain.east);
  \draw[->] (feedbackgain.west) -| (sum.south);
\end{tikzpicture}
\end{center}

\subsection*{(b) $y[n]=\frac{1}{3}y[n-1]+x[n-1]$}

\begin{center}
\begin{tikzpicture}[
  block/.style={draw,minimum width=1.1cm,minimum height=0.8cm},
  sum/.style={draw,circle,minimum size=0.65cm,inner sep=0pt},
  >=stealth,thick]
  \node[block] (inputdelay) at (-2,0) {$\delta[n-1]$};
  \node[sum] (sum) at (0,0) {$+$};
  \node[block] (feedbackgain) at (1.25,-2) {$\frac13\delta[n-1]$};
  \draw[->] (-4,0) node[left] {$x[n]$} -- (inputdelay.west);
  \draw[->] (inputdelay.east) -- (sum.west);
  \draw[->] (sum.east) -- (4,0) node[right] {$y[n]$};
  \fill (2.5,0) circle (1.8pt);
  \draw[->] (2.5,0) |- (feedbackgain.east);
  \draw[->] (feedbackgain.west) -| (sum.south);
\end{tikzpicture}
\end{center}

\section* {2.39 画出微分因果LTI的框图表示}

在初始松弛条件下，将积分合并为一个积分环节。

\subsection* {(a) $y(t)=-\frac{1}{2}\frac{dy(t)}{dt}+4x(t)$}
\[
  \frac{dy(t)}{dt}=8x(t)-2y(t)
  \quad\Longrightarrow\quad
  y(t)=u(t)*\bigl[8x(t)-2y(t)\bigr].
\]
\begin{center}
\begin{tikzpicture}[
  block/.style={draw,minimum width=1.2cm,minimum height=0.8cm},
  sum/.style={draw,circle,minimum size=0.65cm,inner sep=0pt},
  >=stealth,thick]
  \node[block] (inputgain) at (-2,0) {$8\delta(t)$};
  \node[sum] (sum) at (0,0) {$+$};
  \node[block] (integrator) at (2,0) {$\displaystyle\int$};
  \node[block] (feedbackgain) at (1.8,-2) {$-2\delta(t)$};
  \draw[->] (-4,0) node[left] {$x(t)$} -- (inputgain.west);
  \draw[->] (inputgain.east) -- (sum.west);
  \draw[->] (sum.east) -- (integrator.west);
  \draw[->] (integrator.east) -- (4.5,0) node[right] {$y(t)$};
  \fill (3.6,0) circle (1.8pt);
  \draw[->] (3.6,0) |- (feedbackgain.east);
  \draw[->] (feedbackgain.west) -| (sum.south);
\end{tikzpicture}
\end{center}

\subsection* {(b) $\frac{dy(t)}{dt}+3y(t)=x(t)$}
\[
  \frac{dy(t)}{dt}=x(t)-3y(t)
  \quad\Longrightarrow\quad
  y(t)=u(t)*\bigl[x(t)-3y(t)\bigr].
\]
\begin{center}
\begin{tikzpicture}[
  block/.style={draw,minimum width=1.2cm,minimum height=0.8cm},
  sum/.style={draw,circle,minimum size=0.65cm,inner sep=0pt},
  >=stealth,thick]
  \node[sum] (sum) at (0,0) {$+$};
  \node[block] (integrator) at (2,0) {$\displaystyle\int$};
  \node[block] (feedbackgain) at (1.8,-2) {$-3\delta(t)$};
  \draw[->] (-4,0) node[left] {$x(t)$} -- (sum.west);
  \draw[->] (sum.east) -- (integrator.west);
  \draw[->] (integrator.east) -- (4.5,0) node[right] {$y(t)$};
  \fill (3.6,0) circle (1.8pt);
  \draw[->] (3.6,0) |- (feedbackgain.east);
  \draw[->] (feedbackgain.west) -| (sum.south);
\end{tikzpicture}
\end{center}

\newpage
\raggedbottom
\section* {2.40 求系统响应}

\subsection* {(a) 求单位冲激响应}
系统的输入输出关系为
\[
  y(t)=\int_{-\infty}^{t}e^{-(t-\tau)}x(\tau-2)\,d\tau.
\]
令 $x(t)=\delta(t)$，得到

\begin{align*}
  h(t) &=\int_{-\infty}^{t}e^{-(t-\tau)}\delta(\tau-2)\,d\tau \\
       &=\boxed{\begin{cases}
          0, &t<2,\\[0.3em]
          e^{-(t-2)}, &t\ge2.
        \end{cases}}
\end{align*}

\subsection* {(b) 根据给定输入求响应}

\begin{align*}
  y(t) &=x(t)*h(t)
        =\int_{-\infty}^{\infty}x(t-\tau)h(\tau)\,d\tau \\
       &=\int_{2}^{\infty}x(t-\tau)e^{-(\tau-2)}\,d\tau \\
       &=\begin{cases}
          0, &t<1,\\[0.5em]
          \displaystyle\int_{2}^{t+1}e^{-(\tau-2)}\,d\tau,
            &1\le t\le4,\\[0.8em]
          \displaystyle\int_{t-2}^{t+1}e^{-(\tau-2)}\,d\tau,
            &t>4.
        \end{cases}
\end{align*}
计算得
\[
  \boxed{
  y(t)=\begin{cases}
    0, &t<1,\\[0.4em]
    1-e^{1-t}, &1\le t\le4,\\[0.4em]
    e^{-(t-1)}(e^3-1), &t>4.
  \end{cases}}
\]

\newpage
\section* {2.48}

\noindent (a) 正确。取$x(t)=h(-t)$，则$y(0)=\int_{-\infty}^{\infty}x(-\tau)h(\tau)d\tau=\int_{-\infty}^{\infty}(h(\tau))^2d\tau=\infty$

\noindent (b) 不正确。因果 LTI 系统 $y(t)=x(t-1)$ 的逆系统 $x(t)=y(t+1)$ 是非因果的

\noindent (c) $K=0$时显然正确,$k\ne0$时显然不正确:构造$h[n] \equiv \frac{K}{2}$即可。

\noindent (d) 正确。因为$h[n]$有限，显然$\sum_{k-\infty}^{\infty}|h(k)|<\infty$

\noindent (e) 不正确。构造因果LTI$y(t)=\int_{-\infty}^{t}x(\tau)d\tau$，对于有界输入$x(t)\equiv 1$会爆破。

\noindent (f) 不正确。由(b)显然。非因果和因果系统的级联是identity，显然因果。

\noindent (g) 不正确。构造$h(t)=e^{-t}u(t)$,$\int_{-\infty}^{\infty}|h(t)|dt=\int_{0}^{\infty}e^{-t}dt=1$稳定，但是$s(t)=\int_{-\infty}^{t}h(\tau)d\tau=(1-e^{-t})u(t), \int_{-\infty}^{\infty}|s(t)|dt=\int_{0}^{\infty}(1-e^{-t})dt=\infty$

\noindent (h) 正确。因为$h[n]=s[n]-s[n-1]$,$s[n]=0,\forall n<0 \Leftrightarrow h[n]=0,\forall n<0$

\section* {2.51}
\subsection* {(a) $h[n]=(\frac{1}{2})^nu[n],z[n]=nw[n]$系统级联不可交换}

\noindent 取输入信号$x[n] \equiv 1$ \\
对于第一种级联方式：
\[
  w[n]=\sum_{k=-\infty}^{\infty}x[n-k]h[k]=\sum_{k=0}^{\infty}(\frac{1}{2})^k=2 
\]

$$y_1[n]=z[n]=nw[n]=2n$$

\noindent 对于第二种级联方式：
$$z[n]=nx[n]=n$$

\[
  y_2[n]=\sum_{k=-\infty}^{\infty}z[n-k]h[k]=\sum_{k=0}^{\infty}(n-k)(\frac{1}{2})^k=2n-2
\]

\noindent 显然$y_1[n]\ne y_2[n]$，不可交换。

\subsection* {(b) 替换为$z[n]=w[n]+2$重新计算}

\noindent 取输入信号$x[n] \equiv 1$ \\
对于第一种级联方式：
$$w[n]=2$$
$$y_1[n]=w[n]+2=4$$

\noindent 对于第二种级联方式：
$$z[n]=x[n]+2=3$$
\[
  y_2[n]=\sum_{k=-\infty}^{\infty}z[n-k]h[k]=\sum_{k=0}^{\infty}3(\frac{1}{2})^k=6
\]

\noindent 同样有$y_1[n]\ne y_2[n]$不可交换。

\end{document}
